\section{Engineering the proof}
\label{sec:engineering}

This section reports how the result was produced, including what went
wrong. We think the method is reusable, and the failures are the
informative part.

\subsection{Stages and fixed statements}

The work had three stages.

\emph{Stage 0: the semantics.} The specification was written against the
server code at a pinned commit, with the rules citing the files and lines
that implement them, and every modeling choice marked as an assumption or
a decision. There is one deliberate exception: the rule for the commit of
an RPC follows a fix that was not yet merged at the pinned commit (the
first defect of \S\ref{sec:server-bugs}), and is marked as doing so. Three
adversarial reviews checked it against the code. They produced 21 findings;
one was rejected with an argument from the code (the prefix rule of
\rulename{TIMER-FIRE}), and the others changed the specification. The
additions include the fair maximal traces with markers, the admissible
environment and the projection, attempted requests as the input word, and
the late, ordered \kw{fire} action of \Lang{}. The specification ends with a
list of 41 numbered assumptions and decisions.

\emph{Stage 1: the model and the statements.} The Lean model, the
translation, and the two theorems were written first, with the theorems
reduced to inclusion lemmas that ended in \texttt{sorry}: twelve in all.
A review compared the statements with the specification before any proof
started. That review led to the owner decisions L1 to L4 of
\S\ref{sec:language} and to two corrections of the model: snapshots had
loaded every Attribute, where the server loads an AttributeMap instance
only when it is selected; and a runtime-only instruction had been exposed
as a source statement.

\emph{Stage 2: the proofs.} The definitions that the proofs share were
\emph{fixed}: the relation, the invariants, the shape of each step lemma,
and later the liveness definitions. The remaining work was cut into
packages with fixed statements (Table~\ref{tab:packages}), which were
proved in parallel on separate branches and then merged. A package may add
lemmas. It may not change a model file, a fixed definition, or a statement.

\begin{table}[t]
\centering
\small
\begin{tabularx}{\textwidth}{@{}P{1.35cm}LP{5.4cm}@{}}
\toprule
Package & Content & Main statement \\
\midrule
WP1 & Every \Lang{} transition keeps the \Lang{} invariant. &
  \lean{L.Invariant.step} \\
WP2 & \Lang{} to Dex: ticks, requests, timer fires, late emits, the begin of
  an action, the delivery of an effect. &
  \lean{l_tick_simulated}, \lean{l_fire_simulated}, \ldots \\
WP3 & \Lang{} to Dex: the registration of a wait and the completion of an
  action. & \lean{l_thread_register_simulated},
  \lean{l_thread_complete_simulated}, \ldots \\
WP4 & Dex to \Lang{}: every Dex step. & \lean{dex_step_simulated} \\
\midrule
L1 & Block runs flatten; the flattening keeps the trace, the requests, and
  finiteness of progress. & \lean{exists_flattening} \\
L2 & Silent actions lower the silent weight; labels are progress steps
  exactly when the action is one. &
  \lean{silentWeight_lt_of_isSilentKind} \\
L3 & The \Lang{} run of a Dex run: existence, word, trace, marker. &
  \lean{exists_lRun_of_dexRun} \\
L4 & The \Lang{} run of a fair Dex run is fair. & one lemma per assumption \\
L5 & The Dex run of an \Lang{} run: existence, word, trace, marker. &
  \lean{exists_dexRun_of_lRun} \\
L6 & The Dex run of a fair \Lang{} run is fair. & one lemma per assumption \\
\bottomrule
\end{tabularx}
\caption{The work packages. None of them changed a definition of a model
file, a definition of \file{TraceSets.lean}, or a statement of
\file{Theorems.lean}.}
\label{tab:packages}
\end{table}

Fixing the statements first has an obvious benefit: independent provers
cannot drift apart. It also has a cost, which we met at once. A step lemma
quantifies over \emph{every} pair of configurations that satisfies
\lean{IsSimulationState}, not only over reachable pairs. If the fixed
definitions admit an unreachable pair that breaks a step, the lemma is
false, and the prover of that package is the one who finds out.

\subsection{Six gaps, and a fuzz that finds them}
\label{sec:gaps}

That is what happened. The first fixed definitions were not inductive. The
provers found six pairs of configurations that satisfied the definitions,
were not reachable, and falsified a step lemma (Table~\ref{tab:gaps}). Each
gap was closed by one more clause.

\begin{table}[t]
\centering
\small
\begin{tabularx}{\textwidth}{@{}P{2.6cm}LP{4.3cm}@{}}
\toprule
Gap & A pair the first definitions allowed & Closed by \\
\midrule
\kw{WaitFor} task with an effect & A thread that began a wait, whose
  \kw{WaitFor} task holds an effect: the Worker's effect step has no
  \Lang{} step. & \lean{PhaseMatchesStep}: such a task holds no effect. \\
Task effects & After the close, a stale task holds two effects: abandoning
  it drops both, \Lang{} drops one late emit per label. &
  \lean{IsTaskComputed}: a task holds what its handler computes, at most
  one effect. \\
Waiting Step with an invocation & A thread that is not blocked, carried by a
  Step in phase $\mathit{WT}$ with a stored wait: a publication wakes the
  Step and not the thread. & \lean{PhaseMatchesStep} fixes the Step's phase
  in every case. \\
Cancel flag of an unforked child & A flag is set for the next child of a
  live thread, with no pending request: after \kw{fork} the child has a
  flag and no request. & \lean{cancelFlagsAreRequests}. \\
Boot and static mutexes & While \kw{main} boots, the mutex owners and the
  token Channels are unconstrained, and the start Step then publishes the
  tokens. & \lean{mutexesBeforeBoot}. \\
Duplicate timer IDs & The timer list holds one ID twice: the relation reads
  the first binding, the due-timer computation reads both. &
  \lean{TimersAreDistinct}. \\
\bottomrule
\end{tabularx}
\caption{The six gaps in the first fixed definitions. Each row is an
unreachable pair that the definitions accepted and that breaks a step
lemma.}
\label{tab:gaps}
\end{table}

A seventh problem was in the shape of a statement, not in an invariant. The
\Lang{}-to-Dex lemmas first concluded with an existential (``some Dex
actions simulate the label''). The Dex-to-\Lang{} direction could not reuse
them by determinism, because the actions were hidden. The statement was
changed to name them (\lean{simulatingActions},
\S\ref{sec:ltodex}).

Finding such gaps by proving is expensive: each one stops a package, and
the fix touches definitions that other packages rely on. So after the fact
we built the test that would have caught them: an \emph{inductiveness
fuzz} (\file{Examples/Inductiveness.lean}), in the spirit of testing a
property on random inputs before trying to prove it
\cite{claessen2000quickcheck}.
\begin{enumerate}
\item It runs an \Lang{} program and its Flow in lockstep, in both
  directions, on pseudo-random schedules, and samples the reached pairs of
  configurations.
\item It \emph{perturbs} each pair in about 270 ways: one field of one side
  (a thread's phase, counter, timer, locals, or code; the cancel flags; the
  mutex owners; a Step's phase, invocation, wait, or input; a task's phase
  or snapshot; added, dropped, or doubled Steps, tasks, messages, timers,
  Attributes), or matching fields of both sides.
\item It keeps a perturbed pair only when an executable Boolean mirror of
  \lean{IsSimulationState} (\lean{checkState},
  \file{Examples/RelatedCheck.lean}) accepts it. The kept pairs are
  mostly unreachable, which is the point.
\item From every kept pair it takes every enabled \Lang{} label with its
  simulating actions, and every enabled admissible Dex action with its
  \Lang{} labels. The other side must be enabled, the events and clock
  stamps must agree, and the mirror must accept the result.
\end{enumerate}

The result is the honest measure of the method. With the mirror of the
\emph{first} definitions, the run \texttt{inductiveness-fuzz 120 5 6 2}
reports 322 failures, and they cover all six gaps of
Table~\ref{tab:gaps}. With the present definitions the same run, on 13
programs, perturbs 1,898 pairs into 510,426 copies, keeps 189,045 of them,
and reports no failure. A larger run on the finished proof,
\texttt{inductiveness-fuzz 200 4 6 3}, perturbs 5,876 pairs into 1,820,177
copies, again with no failure. The fuzz found no gap beyond the six. It is
evidence, not proof; the work packages prove inductiveness. But the
smaller run takes well under a minute, and had it existed before the statements were fixed,
none of the six gaps would have reached a prover.

Writing the fuzz also taught a lesson about the tooling. The first mirror
of the timer clause checked every binding of the timer list, where the
clause itself reads only the first binding of an ID. A mirror that is
stricter than its clause \emph{hides} gaps: the duplicate-timer gap was
invisible until the mirror was made exact.

\subsection{Testing before proving}
\label{sec:testing}

The same habit was applied throughout. None of these checks is part of the
proof. All of them run when the development is built, and a failure fails
the build.

\emph{Lockstep checks} (\file{Examples/Lockstep.lean},
\file{Examples/LockstepChecks.lean}) run a program and its Flow side by
side on pseudo-random schedules, in both directions, for 13 programs that
together use every construct. After every step the \Lang{} events must
equal the projected Dex events, and the \Lang{} clock stamp must equal the
Dex commit timestamp. Three faults were planted to see whether the build
notices. Dropping the cancel of a blocked timed operation failed. Dropping
the load selection of the waiter list failed. Dropping the timer-number
tie-break \emph{passed}: equal deadlines are rare under random ticks. A
dedicated example with two equal deadlines (\file{Examples/TimerOrder.lean})
was added and fails under that fault.

\emph{Relation checks} (\file{Examples/RelatedCheck.lean}) evaluate the
mirror of \lean{IsSimulationState} in every pair that the lockstep runs
reach, for 14 programs.

\emph{Liveness checks} (\file{Examples/LivenessChecks.lean}) tested the
liveness definitions before the liveness packages started: that
\lean{fairActions} are enabled and end in a settled pair; that the fairness
obligations of the two sides correspond; that every silent action lowers
the silent weight; and that, in a settled pair, \Lang{} has a step other
than a tick or a request exactly when the Flow has one, which is what the
markers \mterm{} and \mdead{} compare. The run
\texttt{inductiveness-fuzz 200 4 6 3 liveness} checks 1,396,813 pairs,
1,205,008 of them settled and 408,933 settled after the close, with no
failure. Each of the four checks was shown to detect a planted fault.

These checks cannot test a statement about an infinite run, and they did
miss one thing. A fixed statement of package L2 was false as first written:
that every silent action lowers the silent weight. With two tasks of one ID
it fails, because an update by ID changes both. The statement now assumes
\lean{TasksAreDistinct}, which every configuration of a run satisfies. The
checks did not find this because they start from pairs that the mirror
accepts, and those have distinct task IDs. The prover found it.

\subsection{What the work found in the real server}
\label{sec:server-bugs}

Writing a rule for every server behavior, with a citation, means reading
every caller of every function the rule mentions. Two defects in the real
server surfaced this way, before any Lean was written.
\begin{itemize}
\item \emph{A cancellation that skipped the graceful-completion check.}
  A Flow that has requested graceful completion closes once no Step is
  active. One path that cancels Steps, cancellation by an RPC, did not run
  that check, so a Flow whose last active Step was canceled this way stayed
  running forever.
\item \emph{Consumed messages lost when a waiting Step was canceled.} A
  Step whose wait had consumed Channel messages, and which was canceled
  before its \kw{Execute} committed, discarded those messages. A review
  finding showed the consequence for a token-based mutex: the token is
  gone and the mutex stays locked.
\end{itemize}
Both have since been fixed in the server. The first is outside the
translation, which uses neither graceful completion nor RPCs. For the
second, the fix changed how the server handles the consumed messages of a
canceled Step. A translated Flow cancels Steps only by closing, and we
argue that the changed handling is unobservable there; but the model
describes the server at the pinned commit, and we have not proved that
the current server refines it (\S\ref{sec:limitations}). The second defect
also shaped \Lang{}: it is why \Lang{}'s cancellation is cooperative and
\kw{lock} is not a cancellation point. The
specification also records where the product documentation said less than
the code does; the plainest case is ``a Step method commits atomically'',
which is true of the writes and not of the read-modify-write
(\S\ref{sec:dex-rules}).

\subsection{Size and effort}
\label{sec:size}

Table~\ref{tab:size} gives the size of the development, counted with
\texttt{wc -l}, comments and blank lines included. The definitions that a
reader must trust are about a tenth of it. The largest part by far is the
\Lang{}-to-Dex step simulation; its largest file, the registration of a
wait, has 1,941 lines.

\begin{table}[t]
\centering
\small
\begin{tabular}{@{}lrrl@{}}
\toprule
Directory & Files & Lines & Content \\
\midrule
\texttt{Core/} & 4 & 1,177 & transition systems, traces, block runs \\
\texttt{Dex/} & 9 & 3,378 & the Dex model, its run invariant, lemmas \\
\texttt{L/} & 11 & 6,033 & the \Lang{} model, its invariant, the lock discipline \\
\texttt{Translation/} & 2 & 433 & the translation, the projection, \Eadm \\
\texttt{Simulation/} & 32 & 15,862 & the relation and the step simulations \\
\texttt{Liveness/} & 16 & 6,435 & the liveness definitions, packages L2 to L6, resting runs \\
\texttt{Examples/} & 13 & 2,692 & kernel-checked runs, witnesses, checks, the fuzz \\
top level & 9 & 2,076 & invariants, finite traces, the theorems \\
\midrule
Total & 96 & 38,086 & \\
\bottomrule
\end{tabular}
\caption{The Lean development at commit \texttt{33e998003}. The
specification is a further 2,005 lines of Markdown.}
\label{tab:size}
\end{table}

The development uses Lean 4.35.0 and no library beyond the Lean core: no
Mathlib and no Batteries. The list and arithmetic facts it needs are in the
core or are proved in the development. Finite maps are association lists,
so that every model function evaluates in the kernel and the examples are
checked by evaluation. A clean build takes a few minutes at most; the two
fuzz runs quoted above took about half a minute and two minutes when we
reran them for this paper.

The proofs were produced by AI coding agents working in parallel on the
work packages with fixed statements, under the author's direction. The
author made all semantic decisions. By their git author dates, the 92
commits under \file{formal/} up to the revision described here fall
between 13:39 and 23:11 on one day: from the first committed draft of the
specification to the complete proof took a little under ten hours. That
figure excludes earlier preparation: planning, the Workflow Patterns
evaluation, and drafting the specification before its first commit.
